% Auto-generated by plot_dp_cbpred_study.py — edit freely
\paragraph{dpPred vs.~dpPred+cbPred and PFQ hit rate.}

For SPEC (DPC3), dpPred alone achieves a geomean IPC speedup of \SI{-0.00}{\percent} over LRU; cbPred adds a further \SI{0.16}{\percent} geomean speedup, for a combined total of \SI{0.16}{\percent}. The PFQ hit rate is moderate at \SI{23.8}{\percent}, consistent with the moderate LLC DoA-on-DoA-page correlation of \SI{35.0}{\percent} seen in the liveness study.
For GAP, dpPred alone achieves a geomean IPC speedup of \SI{0.05}{\percent} over LRU; cbPred adds a further \SI{3.33}{\percent} geomean speedup, for a combined total of \SI{3.38}{\percent}. The PFQ hit rate is moderate at \SI{32.1}{\percent}, consistent with the low LLC DoA-on-DoA-page correlation of \SI{17.1}{\percent} seen in the liveness study.
For Ligra, dpPred alone achieves a geomean IPC speedup of \SI{-0.07}{\percent} over LRU; cbPred adds only \SI{0.07}{\percent} additional speedup, for a combined total of \SI{-0.00}{\percent}. The PFQ hit rate is high at \SI{55.0}{\percent}, consistent with the high LLC DoA-on-DoA-page correlation of \SI{70.7}{\percent} seen in the liveness study.

Across suites, the marginal benefit of cbPred is weakly correlated with the PFQ hit rate ($r=-0.29$), suggesting that other factors beyond PFQ signal quality govern cbPred's marginal contribution on these workloads.
